\documentclass[12pt]{article}

% Modelo de relatório do Trabalho 01, baseado no template de conferências da
% Sociedade Brasileira de Computação (SBC), versão atualizada em 2017.
\usepackage{sbc-template}
\usepackage{graphicx,url}
\usepackage[utf8]{inputenc}
\usepackage[T1]{fontenc}
\usepackage{amsmath,amssymb}
\usepackage{booktabs}
\usepackage{tabularx}
\usepackage{multirow}
\usepackage{float}
\usepackage[hidelinks]{hyperref}

\sloppy
\renewcommand{\tablename}{Tabela}
\renewcommand{\figurename}{Figura}
\renewcommand{\refname}{Referências}

\title{Trabalho 01 - Busca Não Informada, Busca Informada e Busca Local}

\author{Hugo Henrique Marques\inst{1}, João Pedro Silva Siqueira\inst{1}}

\address{Curso de Engenharia de Computação - CEFET-MG, campus Divinópolis\\
Disciplina de Inteligência Artificial - 2026/2
\email{hugohmarques4@gmail.com, joaopedrosilvasiqueira1@gmail.com}}

\begin{document}

\maketitle

\begin{abstract}
This report presents the implementation and comparison of search algorithms in two problems. In a weighted grid, we compare BFS, DFS, Uniform-Cost Search, Greedy Best-First Search and A*, using the Manhattan and Euclidean heuristics. In the 8-Queens problem, we compare Hill Climbing and Random-Restart Hill Climbing with seeds 0 to 99. On the public map, BFS, DFS and Greedy found a path with 15 actions and cost 40, while UCS and A* found the minimum cost of 29. A* with Manhattan expanded 43 states, against 74 for UCS. Hill Climbing solved only 12\% of the cases, while Random Restart solved all of them, with 6.76 restarts on average.
\end{abstract}

\begin{resumo}
Este relatório apresenta a implementação e a comparação de algoritmos de busca em dois problemas. Em um mapa ponderado, comparamos BFS, DFS, Busca de Custo Uniforme, Busca Gulosa e A*, com as heurísticas Manhattan e Euclidiana. No problema das 8 Rainhas, comparamos Hill Climbing e Random-Restart Hill Climbing com as sementes de 0 a 99. No mapa público, BFS, DFS e Gulosa encontraram um caminho com 15 ações e custo 40, enquanto UCS e A* encontraram o custo mínimo de 29. O A* com Manhattan expandiu 43 estados, contra 74 da UCS. O Hill Climbing resolveu apenas 12\% dos casos, e o Random Restart resolveu todos, com média de 6,76 reinícios.
\end{resumo}

\section{Introdução}
O trabalho tem duas partes. Na Parte 1, um agente precisa sair de \texttt{S} e chegar em \texttt{G} em um mapa em grade com obstáculos e terrenos de custos diferentes. Para isso, usamos três buscas não informadas (BFS, DFS e UCS) e duas informadas (Gulosa e A*). Na Parte 2, usamos busca local (Hill Climbing e Random Restart) para posicionar 8 rainhas em um tabuleiro sem ataques. Os algoritmos seguem a definição de \cite{russell2021}.

A análise procura responder principalmente: (i) se o caminho com menos ações também é o mais barato; (ii) qual o efeito das heurísticas no número de estados expandidos; e (iii) por que o Hill Climbing falha e quanto os reinícios ajudam.

\section{Representação e implementação}
\subsection{Parte 1 - Busca não informada e informada}
O estado é a posição \texttt{(linha, coluna)} do agente. O mapa é lido de um arquivo texto, e as posições de \texttt{S} e \texttt{G} são encontradas na leitura, sem dimensões fixas. Os sucessores são gerados na ordem Norte, Leste, Sul e Oeste, ignorando bordas e obstáculos (\texttt{\#}). O custo de uma ação é o custo da célula de destino. Cada busca guarda um dicionário de predecessores e da ação usada, que serve para reconstruir o caminho no final, e um dicionário com o custo $g(n)$.

A BFS usa uma fila (\texttt{deque}) e a DFS usa uma pilha (lista). Na DFS, os sucessores são empilhados em ordem inversa, para que o Norte seja o primeiro a sair da pilha. Nas duas, o estado é marcado como descoberto quando entra na fronteira, usando um conjunto da fronteira e um conjunto de explorados.

UCS, Gulosa e A* usam fila de prioridade (\texttt{heapq}) com as tuplas $(g, \text{ordem})$, $(h, \text{ordem})$ e $(f, h, \text{ordem})$, onde \emph{ordem} é um contador de inserções usado no desempate. Na UCS e no A*, se aparece um caminho mais barato para um estado que ainda está na fronteira, o custo e o predecessor são atualizados e o estado é inserido de novo. A entrada antiga é descartada quando sai da fila e não conta como expansão.

As métricas seguem o enunciado: \emph{gerados} conta todas as inserções na fronteira (incluindo o estado inicial e as reinserções); \emph{expandidos} conta os estados retirados para os quais os sucessores foram gerados (o objetivo não conta); o pico da fronteira é o maior número de estados distintos nela; e o pico de memória é o maior valor de $|\text{fronteira} \cup \text{explorados}|$.

\subsection{Heurísticas}
As duas heurísticas usam a posição do estado e a do objetivo: $h_M(n) = |x_n - x_G| + |y_n - y_G|$ e $h_E(n) = \sqrt{(x_n - x_G)^2 + (y_n - y_G)^2}$. Como o agente só anda em quatro direções, a Manhattan é exatamente o número mínimo de movimentos até o objetivo quando não há obstáculos. A Euclidiana é a distância em linha reta e é sempre menor ou igual à Manhattan.

\subsection{Problema das 8 rainhas}
O tabuleiro é um vetor de 8 posições, em que o índice é a coluna e o valor é a linha da rainha. A função $h(s)$ conta os pares de rainhas que se atacam, isto é, que estão na mesma linha ou na mesma diagonal ($|l_1 - l_2| = |c_1 - c_2|$). Os vizinhos são gerados da coluna 0 até a 7 e, em cada coluna, da linha 0 até a 7, pulando a linha atual, o que dá 56 vizinhos.

O Hill Climbing avalia os 56 vizinhos e escolhe o primeiro com o menor $h$. Ele só se move se esse $h$ for estritamente menor que o atual; caso contrário, para. Assim, quando chega em $h = 0$, ele também para, pois nenhum vizinho pode ser melhor. O Random Restart cria um único \texttt{random.Random(seed)} e usa esse mesmo objeto para gerar todos os estados iniciais, com uma chamada \texttt{randrange(8)} por coluna. Ele faz no máximo 51 tentativas (1 inicial + 50 reinícios) e soma as iterações e os estados avaliados de todas elas.

\section{Metodologia experimental}
Os experimentos foram feitos em Linux com Python 3.12.3, usando apenas a biblioteca padrão. Na Parte 1, usamos o mapa público \texttt{data/mapa\_teste.txt} ($8 \times 16$) com as heurísticas Manhattan e Euclidiana na Gulosa e no A*. Cada configuração foi executada 30 vezes, e o tempo apresentado é a mediana medida com \texttt{time.perf\_counter}, sem contar a leitura do arquivo. O tempo varia um pouco entre execuções e é só uma métrica secundária.

Na Parte 2, os dois métodos foram executados com as sementes de 0 a 99, e o Random Restart usou o limite de 50 reinícios.

O script \texttt{scripts/run\_experiments.py} gera todas as tabelas deste relatório, e os testes automatizados ficam na pasta \texttt{tests}. Os testes cobrem os casos mínimos do enunciado (início ao lado do objetivo, mapa sem solução, mapa com ciclos, rotas com o mesmo número de ações e rota curta e cara contra rota longa e barata), além de testes da função de conflitos, da vizinhança, do Hill Climbing e do Random Restart.

\noindent\textbf{Repositório:} \url{https://github.com/hugnarok/Trabalho-IA} \\
\textbf{Tag da entrega final:} \texttt{entrega-final}.

\section{Resultados}
\subsection{Parte 1 - Busca não informada e informada}
A Tabela~\ref{tab:grid} mostra os resultados no mapa público.

\begin{table}[H]
\centering
\caption{Resultados dos algoritmos de busca no mapa de teste.}
\label{tab:grid}
\resizebox{\textwidth}{!}{%
\begin{tabular}{llrrrrrrr}
\toprule
Algoritmo & Heurística & Passos & Custo & Gerados & Expandidos & Pico fronteira & Memória & Tempo (ms) \\
\midrule
BFS     & -          & 15 & 40 & 66 & 59 & 9  & 66 & 0,073 \\
DFS     & -          & 15 & 40 & 21 & 15 & 6  & 21 & 0,022 \\
UCS     & -          & 29 & 29 & 77 & 74 & 6  & 77 & 0,106 \\
Gulosa  & Manhattan  & 15 & 40 & 21 & 15 & 6  & 21 & 0,023 \\
Gulosa  & Euclidiana & 15 & 40 & 21 & 15 & 6  & 21 & 0,027 \\
A*      & Manhattan  & 29 & 29 & 60 & 43 & 15 & 58 & 0,080 \\
A*      & Euclidiana & 29 & 29 & 72 & 60 & 12 & 72 & 0,100 \\
\bottomrule
\end{tabular}}
\end{table}

A Figura~\ref{fig:caminhos} mostra os dois caminhos encontrados. BFS, DFS e Gulosa seguem em linha reta pela primeira linha e passam pelas 5 células \texttt{M}. UCS e A* com Manhattan dão a volta pela parte de baixo do mapa, passando só por células de custo 1.

\begin{figure}[H]
\centering
\begin{minipage}[t]{0.48\textwidth}
\centering
\small
\begin{verbatim}
S**************G
.#.#####.####...
.#.....#....#...
.#####.#.##.#...
.....#.#....#...
####.#.######...
.....#...R.R....
................
\end{verbatim}
(a) BFS, DFS e Gulosa: custo 40
\end{minipage}
\hfill
\begin{minipage}[t]{0.48\textwidth}
\centering
\small
\begin{verbatim}
S**MMMMM.....**G
.#*#####.####*..
.#*****#....#*..
.#####*#.##.#*..
.....#*#....#*..
####.#*######*..
.....#***R.R**..
........*****...
\end{verbatim}
(b) UCS e A* Manhattan: custo 29
\end{minipage}
\caption{Caminhos no mapa público. \texttt{S} é o início, \texttt{G} é o objetivo e \texttt{*} marca o percurso.}
\label{fig:caminhos}
\end{figure}

\subsection{Parte 2 - Busca local e 8 rainhas}\label{sec:queens}
A Tabela~\ref{tab:queens} resume os experimentos com as sementes de 0 a 99.

\begin{table}[H]
\centering
\caption{Resultados dos métodos de busca local no problema das 8 rainhas.}
\label{tab:queens}
\begin{tabularx}{\textwidth}{Xrr}
\toprule
Métrica & Hill Climbing & Random Restart \\
\midrule
Execuções & 100 & 100 \\
Soluções encontradas & 12 & 100 \\
Taxa de sucesso (\%) & 12 & 100 \\
Média de iterações & 3,28 & 24,72 \\
Mediana de iterações & 3 & 17,5 \\
Média de estados avaliados & 240,68 & 1826,64 \\
Média de reinícios & N/A & 6,76 \\
Máximo de reinícios & N/A & 42 \\
\bottomrule
\end{tabularx}
\end{table}

\textbf{Caso de falha do Hill Climbing (semente 0).} O estado inicial foi \texttt{[6, 6, 0, 4, 7, 6, 4, 7]}, com $h = 9$. O algoritmo fez 4 iterações ($h$: 9, 6, 4, 3, 2), avaliou 281 estados e parou em \texttt{[6, 3, 5, 2, 1, 6, 4, 7]}, com $h = 2$. Os dois conflitos que sobraram são entre as colunas 0 e 5 (mesma linha) e entre as colunas 3 e 4 (diagonal). Entre os 56 vizinhos desse estado, o menor valor de $h$ é 2, e 3 vizinhos têm esse valor. Como nenhum vizinho tem $h$ menor que 2, não existe melhora estrita, e o algoritmo para sem solução.

\section{Discussão}
\begin{enumerate}
\item \textbf{BFS encontrou o menor número de ações? Foi o de menor custo?} Sim, a BFS encontrou 15 ações, que é o mínimo possível, pois \texttt{G} está 15 colunas à direita de \texttt{S}. Mas esse caminho não foi o mais barato: ele passa pelas 5 células \texttt{M} (custo 6 cada) e custa 40, contra 29 da UCS e do A*.

\item \textbf{Por que BFS e UCS podem dar soluções diferentes?} A BFS expande por níveis, ou seja, só olha o número de ações e trata todos os passos como iguais. A UCS expande pelo menor custo acumulado $g(n)$. Quando os custos variam, o caminho com menos ações pode ser caro. No mapa, a UCS preferiu dar a volta (29 ações) para evitar as células \texttt{M}.

\item \textbf{Como a ordem de sucessores influencia a DFS?} A DFS sempre segue o primeiro sucessor disponível até não poder mais. Como \texttt{S} está no canto superior esquerdo, o Norte não existe e o Leste é o primeiro. Por isso a DFS foi direto pela primeira linha até \texttt{G}, com apenas 15 expansões, e achou o mesmo caminho da BFS. Isso aconteceu por causa da ordem, e não por alguma garantia: com o Sul antes do Leste, a DFS desceria pelo mapa e retornaria um caminho diferente.

\item \textbf{Por que a Gulosa expande poucos estados e ainda dá custo alto?} A Gulosa olha só para $h(n)$ e ignora o custo já gasto. Na primeira linha, cada passo para o Leste diminui a distância até \texttt{G}, então ela foi em linha reta, com 15 expansões (o menor número de todos). Mas ela atravessou as células \texttt{M} sem levar o custo em conta, e o resultado custou 40, cerca de 38\% acima do ótimo.

\item \textbf{A* e UCS deram o mesmo custo? E as expansões?} Sim, ambos encontraram custo 29. Mas a UCS expandiu 74 estados, enquanto o A* com Manhattan expandiu 43 (42\% menos) e o A* com Euclidiana expandiu 60. A heurística puxa a busca na direção do objetivo, enquanto a UCS expande em todas as direções pela ordem de custo. Em troca, o pico da fronteira do A* com Manhattan foi maior (15 contra 6).

\item \textbf{Efeito de trocar Manhattan por Euclidiana.} Na Gulosa, não mudou nada: com as duas heurísticas, andar para o Leste na primeira linha é sempre o que mais diminui $h$, então o caminho e as métricas foram iguais. No A*, o custo continuou 29, mas as expansões subiram de 43 para 60. Como $h_E(n) \le h_M(n)$, a Euclidiana é menos informada, e o A* precisa olhar mais estados antes de confirmar o caminho ótimo.

\item \textbf{As heurísticas são admissíveis?} Sim. Com movimentos só em quatro direções, são necessárias pelo menos $|x_n - x_G| + |y_n - y_G|$ ações para chegar ao objetivo. Como o menor custo de passo é 1 (células \texttt{.} e \texttt{G}), o custo real é sempre maior ou igual a $h_M(n)$. E como $h_E(n) \le h_M(n)$, a Euclidiana também nunca superestima. As duas também são consistentes, pois um passo muda a heurística em no máximo 1 e custa pelo menos 1. Isso explica por que o A* achou o custo ótimo com as duas.

\item \textbf{Por que Hill Climbing e Gulosa não são equivalentes?} A Gulosa procura um caminho e guarda uma fronteira com vários estados e um conjunto de explorados. Se um ramo não leva ao objetivo, ela volta e expande outro estado da fronteira. O Hill Climbing guarda só o estado atual, olha apenas os vizinhos imediatos e só aceita melhora estrita. Ele não volta atrás; quando nenhum vizinho é melhor, ele para, mesmo sem solução. Por isso a Gulosa é completa no mapa finito, e o Hill Climbing resolveu só 12\% dos casos.

\item \textbf{Uma falha do Hill Climbing.} Na semente 0, o algoritmo parou com $h = 2$ (Seção~\ref{sec:queens}). Nenhum dos 56 vizinhos tem $h$ menor que 2, então o estado é um mínimo local. Além disso, 3 vizinhos têm o mesmo $h = 2$, ou seja, o estado está em um platô: existiriam movimentos laterais, mas a versão obrigatória não os permite. A subida foi rápida (4 iterações) e parou perto de uma solução, o que foi comum nos experimentos (média de 3,28 iterações).

\item \textbf{Por que os reinícios ajudam, mas não resolvem as limitações?} Cada reinício começa de um ponto aleatório novo. Se uma tentativa resolve cerca de 12\% dos casos, a chance de falhar nas 51 tentativas é de aproximadamente $0{,}88^{51} \approx 0{,}15\%$. Nos experimentos, o Random Restart resolveu as 100 sementes, com média de 6,76 reinícios (perto do valor esperado $0{,}88/0{,}12 \approx 7{,}3$) e máximo de 42. Mas cada tentativa continua sendo um Hill Climbing comum, que para em mínimos locais e platôs; o reinício só troca o ponto de partida. Com um limite de reinícios, não há garantia de sucesso, e o custo aumenta: foram 1826,64 estados avaliados em média, cerca de 7,6 vezes mais que o Hill Climbing.

\item \textbf{Completude e optimalidade.} Neste trabalho o mapa é finito, os custos são positivos (no mínimo 1) e todas as buscas controlam estados repetidos. A Tabela~\ref{tab:prop} resume a comparação com os resultados obtidos.

\begin{table}[H]
\centering
\caption{Propriedades teóricas e resultado observado no mapa público.}
\label{tab:prop}
\begin{tabular}{lccc}
\toprule
Algoritmo & Completo & Ótimo em custo & Custo obtido \\
\midrule
BFS    & Sim & Não (ótimo em número de ações) & 40 \\
DFS    & Sim (mapa finito) & Não & 40 \\
UCS    & Sim & Sim & 29 \\
Gulosa & Sim (mapa finito) & Não & 40 \\
A*     & Sim & Sim (heurística admissível) & 29 \\
\bottomrule
\end{tabular}
\end{table}

Os resultados batem com a teoria: só UCS e A* garantem o menor custo e foram os únicos a achar 29. BFS, DFS e Gulosa acharam o caminho mais curto em ações, mas mais caro. O Hill Climbing, por outro lado, não é completo, como mostram os 12\% de sucesso.
\end{enumerate}

\section{Conclusão}
No mapa ponderado, o menor número de ações não foi o menor custo: BFS, DFS e Gulosa chegaram com custo 40, e só UCS e A* chegaram ao custo ótimo de 29. O A* com Manhattan foi a melhor opção, pois achou o ótimo com 43 expansões, contra 74 da UCS. A Euclidiana também levou ao ótimo, mas com mais expansões. Nas 8 Rainhas, o Hill Climbing é rápido, mas resolveu só 12\% dos casos por parar em mínimos locais e platôs. Os reinícios aleatórios levaram o sucesso a 100\% nas sementes testadas, ao custo de cerca de 7,6 vezes mais estados avaliados.

\section*{Contribuições dos autores}
\textbf{Hugo Henrique Marques:} estrutura inicial do projeto, heurísticas, função de reconstrução do caminho, BFS, DFS e Busca Gulosa, primeira versão do Hill Climbing e do Random Restart, testes mínimos e script de experimentos da Parte 1, README da Parte 1 e saídas do terminal.

\noindent\textbf{João Pedro Silva Siqueira:} UCS e A*, correção das métricas em caso de falha e do desempenho das buscas, correção da função de conflitos das 8 Rainhas, testes e experimentos da Parte 2, README final e relatório.

\section*{Declaração de uso de IA}

Ferramentas de Inteligência Artificial foram utilizadas para ajudar a fazer os testes, corrigir alguns erros, identificar erros ortográficos e auxiliar na revisão.

\begin{thebibliography}{99}

\bibitem[Russell and Norvig 2021]{russell2021}
Russell, S. and Norvig, P. (2021). \textit{Artificial Intelligence: A Modern Approach}. Pearson, 4th edition.

\end{thebibliography}

\end{document}
